\section{Evaluation}
\label{sec:evaluation}

\paragraph{Framing.} Success rate against a single-body external baseline (PHC/PULSE/ExBody2/GMT)
cannot demonstrate this paper's actual claim — that one policy learns the \emph{same motion}
across a wide body-shape distribution — since those methods only ever evaluate one body.
World-space error ($\mathrm{gt\_error}$, meters) is not a fair axis even for \emph{internal}
cross-shape comparison: a taller body produces larger absolute position deviations for the same
relative motor error, so raw meters bake in a body-scale confound before measuring skill at all.
DOF-angle-space error ($\mathrm{dp\_error}$, radians) is exactly shape-invariant by construction
(same joint target regardless of body) and is therefore the primary lens for the paper's core
claim, not a secondary metric introduced only for the reward-design work of
Section~\ref{sec:exp:reward-source}. Formally,
\begin{equation}
  \mathrm{dp\_error} = \operatorname{mean}_{d}\big| q_{\text{dof}}[d] - q_{\text{dof,ref}}[d] \big|,
  \qquad
  \mathrm{gt\_error} = \operatorname{mean}_{b}\big\| p[b] - p_{\text{ref}}[b] \big\|_2,
  \label{eq:dp-gt-error}
\end{equation}
where $d$ ranges over actuated DOFs (radians) and $b$ ranges over rigid bodies (meters).

\paragraph{Scope note: $\mathrm{dp\_error}$ measures joint rotation, not root trajectory.}
$q_{\text{dof}}$ is the actuated joint-angle vector, local to each body's own kinematic tree —
hinge/spherical joint values relative to their parent link. It excludes the root free-joint
(pelvis translation and orientation in the world), which is tracked separately via
$\mathrm{gt\_error}$/$\mathrm{gr\_error}$ and is exactly the shape-dependent signal deliberately
excluded from the headline metric. The shape-invariance claim built on $\mathrm{dp\_error}$ is
therefore scoped to \textbf{limb articulation} — ``the same joint rotates the same way regardless
of body'' — not to whole-body world-space behavior. A tall and a short body executing the same
clip show identical $\mathrm{dp\_error}$ trajectories even though their root paths differ in
absolute distance (shorter legs, shorter stride, same joint angles); that is expected and outside
what this metric measures. A root-level shape-adaptation claim (e.g., stride length scaling
correctly with leg length) would require a separate, explicitly shape-normalized root metric
(e.g., root displacement normalized by height/leg length) — raw $\mathrm{gt\_error}$ cannot
support that claim either, for the same body-scale-confound reason.

\subsection{Cross-Shape Consistency per Clip}
For each clip, we report the spread (std/IQR) of $\mathrm{dp\_error}$ across the 128 training
shapes. A tight spread demonstrates shape-independent skill execution; a wide spread would
undercut the claim.

\subsection{Shape-Failure Correlation}
Computed across multiple runs (full detail in Section~\ref{sec:exp:failure}): Pearson
$r \approx -0.09$ to $0.003$ between shape extremity (mass/height) and failure — effectively zero
relationship. This is the headline evidence that body shape does not predict success, i.e.\ the
same motion is being learned regardless of body.

\subsection{Held-Out Shape Generalization}
Extends the consistency claim past the 128 training shapes to interpolated/extrapolated unseen
shapes (experiment E7, Section~\ref{sec:exp:cross-shape}) — the strongest version of this
evaluation. \todo{Currently blocked on the \texttt{infer.py} pipeline bug.}

\subsection{External Comparison (Footnote)}
Neutral-shape success rate against single-body baselines (PHC-style, $\beta = 0$) is retained only
as a minor sanity-check — confirming neutral-shape quality was not sacrificed for shape
generalization — not as a headline result, since a single-body number structurally cannot speak to
the shape-generalization claim this section exists to make.
